\section{ZNALEZISKA}

Ogólna zasada:

W pewnych komnatach oddział może się natknąć na osobliwe przedmioty, takie jak fontanny, posągi, meble, ołtarze, lustra, dzieła sztuki, schody czy drzwi-pułapki.
Po pokonaniu wszystkich potworów, które ewentualnie znajdują się w takiej komnacie oddział może (ale nie musi) zbadać znaleziony przedmiot.

Sposób postępowania:

Do zbadania znaleziska musi być wydelegowana spośród członków oddziału jedna postać.
W większości wypadków ponosi ona jako jedyna wszelkie konsekwencje, zarówno negatywne, jak i pozytywne, oddziaływania znalezionego przedmiotu.
Gracz, którego postać została wydelegowana rzuca 1K6 i w tabeli 13.9 ,,Znaleziska'' sprawdza na przecięciu kolumny, odpowiadającej typowi znaleziska oraz rzędu, odpowiadającego liczbie wyrzuconych oczek dokładną nazwę tego, co zostało znalezione.
Opisy poszczególnych znalezisk umieszczone zostały w punktach 13.1 do 13.8.

Zasady szczegółowe:

[13.1] Fontanna.

W komnacie stoi duża, wykonana z brązu fontanna.
Zbadanie fontanny polega na wypiciu pewnej ilości tryskającej z niej cieczy.
Oddział może natknąć się na 6 rodzajów fontann:

Trucizna~-- Ciecz jest trucizną.
Badająca postać natychmiast odnosi 1K3 ran.

Mikstura~-- Ciecz jest magiczną miksturą.
Należy w tabeli 14.9 ,,Magiczne przedmioty'' sprawdzić, rzucając 1K6, jej rodzaj (patrz 14.5) oraz natychmiast wprowadzić w życie efekty wynikające z wypicia tego rodzaju mikstury.
Postać, przeprowadzająca badanie może wziąć ze sobą jedną porcję mikstury (tzn. ilość wystarczającą na jednokrone jej użycie w przyszłości).

Alkohol~-- Postać, przeprowadzająca badanie wypiła niezwykle mocny alkohol i jest teraz pijana.
Jej zręczność ulega do końca przygody zmniejszeniu o 2 punkty.

Brylant~-- Ciecz jest wodą, ale w niej postać, badająca fontanna znajduje duży brylant.
Jego wartość należy określić według tabeli 14.9 ,,Kosztowności'', rzucając 2K6 i dodając 2 punkty do liczby wyrzuconych oczek.

Woda~-- Ciecz jest wodą, nie posiadającą żadnych niezwykłych właściwości.

Krew~-- Ciecz jest krwią, a jej wypicie powoduje chorobę.
Zręczność postaci, przeprowadzającej badanie jest zmniejszona o 1 punkt do końca bieżącej przygody.

[13.2] Posąg.

W komnacie stoi naturalnej wielkości posąg, wykonany ze wspaniale obrobionego alabastru.
Badająca postać musi podejść do niego, by mu się dokładnie przyjrzeć.
Może być 6 rodzajów posągów:

Meduza~-- Posąg przedstawia Meduzę, która w czasie, gdy postać jej się przygląda ożywa, spoglądając na przeprowadzającą badanie postać.
Postać musi próbować odeprzeć czar, a jeżeli próba będzie nieudana zostaje zamieniona w kamień.
Reszta oddziału postępuje tak, jak przy normalnym spotkaniu z potworem (patrz 15.3).

Brylanty~-- Posąg przedstawia wielkiego barana, w którego oczodołach tkwią dwa brylanty.
Przeprowadzająca badanie postać może je zabrać (określając ich wartość przy pomocy tabeli 14.9 ,,Kosztowności'' i rzutu 2K6+2).

Medalion~-- Na szyi posągu wisi medalion, który postać, przeprowadzająca badanie może zabrać.
Jego rodzaj ustalany jest przy pomocy tabeli 14.9 ,,Magiczne przedmioty''.

Demon~-- Posąg przedstawia Demona.
Postać, przeprowadzająca badanie musi rzucić 1K6 i w tabeli 13.9 ,,Znaleziska'' w kolumnie ,,ołtarz'' odnaleźć jego imię.
Posąg jest traktowany jak ołtarz, poświęcony temu Demonowi (patrz 13.5).

Talizman~-- Na szyi posągu wisi niewielki talizman, który może być zabrany przez postać, przeprowadzającą badanie.
Jego rodzaj jest ustalany przy pomocy tabeli 14.9 ,,Magiczne przedmioty''.

Bezimienny~-- Posąg przedstawia odwiecznego i bezlitosnego wroga wszystkiego co żywe~-- Bezimiennego.
Postać, przeprowadzająca badanie musi natychmiast podjąć próbę odparcia czaru.
Jeżeli zakończy się ona pomyślnie~-- nic się nie stanie, jednak jeżeli się nie powiedzie postać staje się narzędziem w rękach Bezimiennego i musi natychmiast zaatakować pozostałych członków oddziału (ustawiana jest naprzeciw środkowej postaci w pierwszym szeregu szyku oddziału).
Rozpoczyna się walka, która może zakończyć się tylko śmiercią zaczarowanej postaci lub członków oddziału, bądź też uwieńczonym powodzeniem rzuceniem na postać czaru ,,uwolnienie'' (postać, znajdująca się we władzy Bezimiennego musi próbować odeprzeć ten czar).

[13.3] Drzwi~-- pułapka.

W podłodze, w centrum komnaty znajdują się drzwi~-- pułapka.
Badanie polega na ich otworzeniu.
Oddział może się natknąć na cztery rodzaje takich drzwi:

Pułapka~-- W drzwiach ukryta jest ,,normalna'' pułapka.
Przeprowadzająca badanie postać musi sprawdzić jej rodzaj w tabeli 7.1 ,,Pułapki''.
Jeżeli pułapka zostanie szczęśliwie unieszkodliwiona lub zadziała, a jej efekty nie zabija postaci, może ona zabrać skarb typu J, który był ukryty pod drzwiami~-- pułapką.

Pokój~-- Przeprowadzająca badanie postać spada do znajdującego się pod podłogą (ale nie na niższym poziomie) pokoju.
Postać musi ustalić czy znajduje się w nim potwór (normalną metodą), a jeśli tak~-- zachowywać się jak przy normalnym spotkaniu.
W pokoju znajduje się skarb typu J.
Po uporaniu się z potworem (jeśli był) postać może zabrać ten skarb i wyjść z powrotem.
Pozostali członkowie oddziału nie mogą jej w żaden sposób pomagać.

Zapadnia~-- Podobnie jak wyżej, z tym że pod podłogą znajdują się kronki, w liczbie równej 1K3.
Postać nie może wyjść dopóki się z nimi nie upora.

Czarne Wrota~-- Postać, przeprowadzająca badanie spada na pochylnię, która wiedzie wprost ku Czarnym Wrotom.
Postać jest eliminowana z gry do czasu zniszczenia Czarnych Wrót (patrz 16.0), potem zaś jest uwolniona z ich mocy i może wrócić do gry.

[13.4] Mebel.

W komnacie stoi magiczny mebel, który może zostać zbadany przez dokładne oględziny.
Oddział może się natknąć na 6 rodzajów mebli:

Trumna~-- W centrum komnaty stoi trumna, z której podnosi się Wampir (patrz 15.5).
Postać, przeprowadzająca badanie musi natychmiast podjąć z nim walkę (nie może negocjować, ani usiłować się wykupić).
Po pierwszej rundzie do walki mogą się włączyć pozostali członkowie oddziału.

Regał~-- Postać, przeprowadzająca badanie rzuca 1K6~-- liczba oczek od 1 do 3 oznacza, że regał przewraca się na nią, zadając 1K3 ran, zaś liczba oczek od 4 do 6 pozwala postaci na znalezienie księgi, z której może się ona nauczyć czaru ,,wskrzeszenie'' (patrz 10.8).

Biurko~-- Pod blatem biurka założona jest pułapka.
Jeżeli postać, przeprowadzająca badanie chce wziąć leżący w szufladzie medalion musi się z tą pułapką uporać (unieszkodliwić bądź wziąć na siebie skutki jej zadziałania~-- patrz 7.0), a następnie przy pomocy tabeli 14.9 ,,magiczne przedmioty'' ustalić rodzaj medalionu.

Łóżko~-- W komnacie stoi magiczne łoże, które może wyleczyć postać, przeprowadzającą badanie (tylko ją) z 1K3 ran.

Klawikord~-- W centrum komnaty stoi klawikord, który sam gra.
Postać, przeprowadzająca badanie musi podjąć próbę odparcia czaru~-- jeżeli będzie ona nieudana klawikord swoją muzyką zmusza postać do włożenia do jego wnętrza połowy posiadanych skarbów.
Ich odzyskanie jest możliwe, ale wiąże się z odniesieniem przy próbie otworzenia klawikordu 2 ran.

Lustro~-- patrz 13.7

[13.5] Ołtarz

W komnacie stoi ołtarz, poświęcony jednemu z sześciu Demonów, którym oddają cześć mieszkańcy labiryntu.
Postać, przeprowadzająca badanie musi próbować odeprzeć czar.
Jeśli się jej to uda ołtarz nagradza ją, jeśli nie~-- rzuca przekleństwo.
Skutki zarówno jednego, jak i drugiego są następujące:

Alloces~-- Jeżeli próba odparcia czaru była udana postać może w czasie najbliższej walki rzucać wszystkie czary typu ,,W'' za darmo, nie odnosząc ran, stanowiących ich koszt.
Jeżeli próba nie powiodła się zręczność postaci jest obniżona o 1 punkt do końca przygody.

Vassago~-- Jeżeli próba odparcia czaru była udana postać, przeprowadzająca badanie zyskuje zdolności ,,pułapki'' równe +3 (lub jeżeli dotychczasowe zdolności ,,pułapki'' ulegają zwiększeniu o 3 punkty, jednak nie mogą być wyższe niż 5 punktów).
Jeżeli próba się nie uda postać traci swe zdolności ,,pułapki'' (jeśli je ma) na stałe, do końca całej gry.

Avnas~-- Jeżeli próba odparcia czaru była udana postać, przeprowadzająca badanie zostaje obdarowywana jest możliwością rzucania czaru ,,błyskawica'' kosztem 1 rany, niezależnie od tego, czy w ogóle posiadała dotychczas umiejętność rzucania tego czaru.
Jeżeli próba się nie uda czar ,,błyskawica'' rzucany jest na postać, przeprowadzająca badanie.

Malthus~-- Postać, która pomyślnie przejdzie próbę odparcia czaru zyskuje umiejętność rzucania czaru ,,Gniew Boży''.
Jeżeli próba się nie powiedzie ołtarz atakuje ją za pomocą czaru ,,Gniew Boży''.

Leraje~-- Postać, która pomyślnie przeszła próbę odparcia czaru zyskuje biegłość w posługiwaniu się łukiem na wysokości +3 (bez względu na to czy posiada łuk, czy nie) lub jej dotychczasowa biegłość jest zwiększana o 3 punkty.
Jeżeli próba się nie powiedzie postać jest atakowana przez 3 magiczne strzały~-- rezultat ataku musi być natychmiast sprawdzony w kolumnie ,,łuk'' tabeli 9.9 ,,Rezultat walki''.

Asmoday~-- Jeżeli próba odparcia czaru była pomyślna zręczność postaci, przeprowadzającej badanie zostaje zwiększona o 3 punkty oraz zyskuje ona umiejętność rzucania czaru ,,błyskawica'' kosztem jednej tylko rany (bez względu na to, czy dotychczas mogła się tym czarem posługiwać czy nie).
Jeżeli próba się nie powiedzie każdy ze składników potencjału magicznego postaci ulega redukcji o 1.

[13.6] Dzieło sztuki.

W komnacie znajduje się dzieło sztuki o dużej wartości artystycznej.
Oddział może się natknąć na 6 rodzajów dzieł sztuki:

Gobelin~-- Na ścianie wisi misternie tkany gobelin, wykonany w pracowniach Elfów.
Postać, przeprowadzająca badanie może go zabrać, musi jednak przedtem wyrzucić jedną ze swych broni.
Wartość gobelinu określa się według tabeli 14.9 ,,Kosztowności'', dodając 4 do liczby wyrzuconych oczek.

Obraz~-- Na ścianie wisi obraz, przedstawiający grupę postaci.
Każdy z członków oddziału musi natychmiast rzucać 1K6.
Jeżeli któryś z nich wyrzuci 6 oczek oznacza to, że jest on jedną z postaci, przedstawionych na obrazie.
Świadomość tego działa jak przekleństwo i przedstawiona na obrazie postać odnosi 1K3 ran.

Rzeźba~-- W komnacie znajduje się duża rzeźba.
Należy ją traktować jak posąg i ustalić jej nazwę oraz właściwości, posługując się kolumną ,,posąg'' tabeli 13.9 ,,Znaleziska''.

Kryształ~-- Na stole w centrum komnaty stoi sześcian ze szkła kryształowego o pięknie rżniętych ścianach, przedstawiających sceny z życia labiryntu.
Jeżeli postać, przeprowadzająca badanie zdecyduje się ten kryształ zabrać musi rzucić 1K6~-- od 1 do 3 oczek oznacza, że kryształ ma własności talizmanu, od 4 do 6 oczek~-- kryształ ma własności medalionu.
Należy posłużyć się odpowiednią kolumną tabeli 13.9 ,,Znaleziska'' w celu ustalenia własności jednego bądź drugiego.

Ikona~-- Na ścianie wisi ikona, na której przedstawiony jest jeden z 6 Demonów, zamieszkujących labirynt.
Należy, posługując się kolumną ,,ołtarz'' tabeli 13.9 ,,Znaleziska'' ustalić imię tego Demona.

Manuskrypt~-- Na stole w centrum komnaty leży manuskrypt, którego przeczytanie pozwoli postaci, przeprowadzającej badanie na opanowanie sztuki rzucania czaru ,,Gniew Boży'' (patrz 10.8).

[13.7] Lustro.

Znalezione lustro może wyjawić gdzie znajdują się Czarne Wrota.
Od pierwszego lustra, na które oddział się natknie postacie dowiedzą się na którym poziomie labiryntu Czarne Wrota się znajdują.
Procedura postępowania polega w tym wypadku na rzuceniu 1K6 i odczytaniu w tabeli 13.9 ,,Lustro'', w kolumnie poziom numeru poziomu.
Jeżeli odnaleziony w ten sposób poziom jest tym samym, na którym oddział aktualnie przebywa wykonywany jest drugi rzut 1K6.
Z kolumny ,,segmenty'' tej samej tabeli odczytywana jest liczba niezbadanych dotychczas segmentów, dzielących miejsce, w którym przebywa oddział od miejsca, w którym znajdują się Czarne Wrota.

Jeżeli poziom, wskazany przez pierwszy rzut jest inny od tego, na którym znalezione zostało pierwsze lustro drugi rzut nie jest wykonywany.
Zamiast tego oddział musi znaleźć na wskazanym przez pierwszy rzut poziomie drugie lustro i dopiero wtedy rzucić drugi raz oraz odczytać liczbę niezbadanych segmentów, dzielącą go od Czarnych Wrót.

Przykład:

Oddział odnajduje lustro na pierwszym poziomie i rzuca kostką.
Liczba oczek wskazuje, że Czarne Wrota znajdują się na poziomie trzecim.
Oddział szuka schodów, po czym schodzi dwa poziomy niżej.
Tam odnajduje drugie lustro i rzuca kostką drugi raz.
Liczba oczek wynosi 4.
Oznacza to, że Czarne Wrota znajdują się w szóstym niezbadanym dotychczas segmencie, licząc następny, do którego oddział wejdzie~-- jako pierwszy.
Po przejściu pięciu segmentów (bez względu na kierunek) w sposób normalny oddział znajduje Czarne Wrota w szóstym, tzn. żeton, który je reprezentuje kładziony jest na tym segmencie.

Uwaga! Ważne jest jedynie pierwsze znalezione lustro i ewentualnie to, na które oddział natknie się na poziomie, na którym według pierwszego lustra mają się znajdować Czarne Wrota~-- inne mogą zostać zignorowane.

[13.8] Schody~-- patrz 6.6.

[13.9] Tabele ,,Znaleziska'', ,,Lustro''~-- patrz plansza z tabelami.
